\documentclass[letterpaper]{article}
\usepackage[submission]{aaai2027}
\usepackage[hyphens]{url}
\usepackage{graphicx}
\urlstyle{rm}
\def\UrlFont{\rm}
\usepackage{natbib}
\usepackage{caption}
\frenchspacing
\usepackage{booktabs}
\usepackage{algorithm}
\usepackage{algorithmic}

\pdfinfo{
/TemplateVersion (2027.1)
}

\setcounter{secnumdepth}{0}

\newcommand{\method}{G-Wrist}
\newcommand{\todo}[1]{\textbf{TODO-#1}}

\title{Geometry-Anchor Guided Wrist-View Video Generation for Robot Manipulation}
\author{Anonymous Submission}
\affiliations{}

\begin{document}

\maketitle

\begin{abstract}
Wrist-view observations are a bottleneck for robot policy learning: they reveal gripper-object contact and near-field occlusion, yet are much less available than fixed external-camera videos in many robot datasets and deployments.
We study wrist-view video synthesis from two synchronized external manipulation videos, robot state/action, and camera geometry, a setting made difficult by large exo-to-wrist viewpoint changes, rapid end-effector motion, and self-occlusion.
Unlike anchor-to-wrist pipelines that rely on reconstructed target-view projections as dense generation conditions, we introduce \method, a geometry-anchor guided diffusion model that uses sparse LagerNVS wrist anchors only as slot-aligned inpainting evidence, fuses external-view latents through temporal source memory, and initializes denoising with an action/view-aware 3D scene-token noise prior.
This design treats geometry as sparse evidence for a robot-conditioned video generator rather than as a complete renderer, and evaluates the generated view by downstream policy utility.
On the current DROID-style visual protocol, \method{} outperforms the evaluated close baselines, reducing FVD from 103.93 for WristWorld to 57.92 and LPIPS from 0.3981 to 0.2174.
On RLBench, generated wrist observations improve pi0.5 replacement training from 38.96\% to 43.82\% and nearly match real-wrist augmentation, reaching 50.18\% versus 50.40\%.
These results indicate that sparse geometry-guided diffusion can convert abundant external videos into policy-useful wrist observations.
\end{abstract}

\section{Introduction}

Robot manipulation policies increasingly rely on visual observations to handle contact-rich interaction, occlusion, and object-state change.
Vision-language-action models and large robot datasets have expanded visuomotor learning \cite{brohan2023rt2,kim2024openvla,khazatsky2025droid}, while robot video prediction and embodied world models suggest that generated observations can support policy learning \cite{hu2025vpp,zhen2025tesseract,jang2025dreamgen}.
In these settings, viewpoint matters.
External cameras capture the workspace layout and object motion, but wrist cameras often provide the most direct evidence of gripper-object contact, fine alignment, and near-field occlusion.
Unfortunately, wrist-camera streams are expensive to collect consistently, may be missing in legacy datasets, and can fail or be unavailable during deployment.

This paper studies a targeted observation-completion problem: given two synchronized external-view manipulation videos, robot state/action, and camera geometry, synthesize the corresponding wrist-mounted camera video.
The problem is related to exocentric-to-egocentric video generation \cite{liu2024exo2ego,xu2025egoexogen,park2026egoworld} and recent wrist-view generation \cite{qian2025wristworld}, but robot manipulation imposes a sharper geometric and control-oriented requirement.
The generated wrist video should not only look plausible; it should preserve the near-field evidence needed by downstream policies.
A visually appealing but contact-inconsistent wrist video can hurt a policy even if it scores well on generic perceptual metrics.

The technical difficulty comes from the mismatch between available observations and the desired target view.
A video diffusion model can synthesize temporally coherent frames, but without target-view evidence it may drift under large exo-to-wrist viewpoint changes.
Neural and geometric novel-view synthesis methods can provide sparse target-view anchors from external views \cite{wang2024dust3r,leroy2024mastr,wang2025vggt,szymanowicz2026lagernvs}, but sparse anchors alone do not define a complete manipulation video across a whole episode.
We therefore treat geometry as structured evidence for diffusion rather than as a complete renderer.

We introduce \method, a geometry-anchor guided wrist-view generator built on a WAN-style latent video diffusion backbone \cite{wan2025wan}.
LagerNVS-derived wrist anchors are inserted at the head, tail, and selected keyframes and encoded through the inpainting branch as a slot-aligned mask-plus-latent condition.
External-view latents are projected into temporal source memory so that the target wrist frames attend to nearby source evidence.
Robot state/action and camera tokens provide dynamic and viewpoint conditioning, while scene tokens drive an action/view-aware 3D noise initialization for the target wrist latent.
The current implementation uses scene tokens through active sidecar, noise-prior, and alignment paths; it does not claim active scene-token injection into every diffusion-transformer block.

Our contributions are:
\begin{itemize}
    \item We integrate sparse LagerNVS wrist anchors into a wrist-view video diffusion model through slot-aligned inpainting conditioning, using geometry as target-view evidence rather than a complete renderer.
    \item We introduce an action/view-aware 3D noise initialization that uses scene tokens, camera/view information, and robot state/action to bias target wrist diffusion while preserving scheduler-compatible noise statistics.
    \item We evaluate generated wrist observations as downstream policy observations on RLBench, covering replacement and augmentation settings instead of relying only on visual metrics.
\end{itemize}

Our current evidence supports three bounded claims.
First, under the current DROID-style visual protocol, \method{} improves over the evaluated close baselines, including WristWorld and Exo2Ego-V, in PSNR, SSIM, LPIPS, FID, and FVD.
Second, ablations show progressive visual gains from LagerNVS anchors, temporal keyframe anchors, scene tokens, and the 3D noise prior.
Third, RLBench policy experiments indicate that generated wrist observations recover much of the real-wrist benefit in the recorded settings and nearly match real-wrist augmentation for pi0.5.
Before final submission, these claims require visual-protocol unification, confidence intervals, downstream component ablations, focused novelty comparison, and code-claim audit.

\section{Related Work}

\textbf{Exo-to-ego and wrist-view video generation.}
Cross-view video generation translates external observations into egocentric videos, typically for human activities or daily-life environments \cite{liu2024exo2ego,xu2025egoexogen,park2026egoworld}.
WristWorld brings this broad question closer to robot manipulation by generating wrist views through 4D world modeling \cite{qian2025wristworld}.
These works motivate the importance of view transfer.
Our draft positions the contribution more narrowly: sparse target-view geometry anchors and action/view-aware diffusion initialization synthesize policy-useful wrist observations from synchronized robot external views.
\todo{LIT}: finalize the closest-work novelty delta against WristWorld, Exo2Ego-V, EgoExo-Gen, EgoWorld/EgoX, and recent exo-to-ego video generation papers.

\textbf{Geometry and novel-view priors.}
Feed-forward 3D reconstruction and geometry-aware view synthesis provide camera, point-map, and scene-token priors from sparse views \cite{wang2024dust3r,leroy2024mastr,wang2025vggt}.
LagerNVS provides a neural latent-geometry path for novel-view synthesis \cite{szymanowicz2026lagernvs}.
Our method does not treat these tools as the final video generator.
Instead, it uses sparse wrist anchors and camera-aware scene tokens as evidence for a temporal diffusion model, allowing the generator to fill contact-rich dynamics between sparse target-view observations.

\textbf{Video diffusion and robot world models.}
Latent video diffusion and diffusion-transformer video generators provide strong temporal synthesis backbones \cite{blattmann2023svd,wan2025wan}.
In robotics, video prediction policies and embodied world models investigate whether predicted observations or generated trajectories can improve control and generalization \cite{hu2025vpp,zhen2025tesseract,jang2025dreamgen}.
Our work follows this policy-grounded direction, but focuses on missing-wrist observation synthesis rather than general future-video generation.
The generated wrist video is evaluated both as pixels and as a training observation for downstream policies.

\textbf{Robot datasets and policies.}
Large-scale robot datasets and VLA models have made multi-task robot learning increasingly practical \cite{khazatsky2025droid,brohan2023rt2,kim2024openvla}.
RLBench provides a controlled manipulation benchmark for policy evaluation \cite{james2019rlbench}, and action-generative policies such as Diffusion Policy show the value of expressive visuomotor output distributions \cite{chi2024diffusionpolicy}.
In this draft, DROID-style data supports visual wrist-view generation evaluation and RLBench supports downstream policy evaluation.
DROID downstream policy results and real-robot results are not currently claimed.

\section{Method}

\subsection{Overview}

\method{} is a target-view stage-2 generator built on a WAN-style latent video diffusion backbone.
Figure~\ref{fig:pipeline} sketches the intended data path.
Two external videos are encoded into VAE latents and projected into temporal source memory.
LagerNVS provides sparse wrist-view anchors at selected target slots, and the inpainting branch encodes these anchors as mask-plus-latent condition $y$.
Robot state/action and camera tokens provide dynamic and viewpoint conditions.
Scene tokens feed a structured 3D noise prior that replaces independent Gaussian initialization with a normalized scene/action/view-aware mixture.

\begin{figure}[t]
\centering
\fbox{\rule{0pt}{0.9in}\rule{0.92\linewidth}{0pt}}
\caption{Method overview placeholder. The final figure should show two external videos, LagerNVS wrist anchors, the slot-aligned inpainting branch, source temporal memory, action/camera conditioning, and the 3D noise initialization path. \todo{FIG}}
\label{fig:pipeline}
\end{figure}

\subsection{Problem Formulation}

Let $X^{(1)}_{1:T}$ and $X^{(2)}_{1:T}$ denote two synchronized external-view RGB videos of a robot manipulation episode, and let $a_{1:T}$ denote the robot action or state sequence.
The target is a wrist-view video $Y_{1:T}$ from a camera mounted near the end effector.
Camera intrinsics and extrinsics, when available, are represented as camera tokens $c$.
The generator learns a conditional distribution
\[
    p_{\theta}(Y_{1:T}\mid X^{(1)}_{1:T}, X^{(2)}_{1:T}, a_{1:T}, c).
\]
In stage-2 training, the model predicts target wrist latents only, while external-view latents serve as conditioning memory.

Let $E_{\mathrm{vae}}$ and $D_{\mathrm{vae}}$ denote the video VAE encoder and decoder.
The target latent is $z_0=E_{\mathrm{vae}}(Y_{1:T})$.
Source-view latents $s^{(1)}$ and $s^{(2)}$ are encoded from $X^{(1)}_{1:T}$ and $X^{(2)}_{1:T}$ and are not part of the stage-2 reconstruction loss.
The conditioning problem is therefore asymmetric: source views provide memory, sparse wrist anchors provide target-view evidence, and the model denoises only the wrist latent.

\subsection{Sparse Geometry Anchors}

The first target-view evidence stream is a sparse set of synthesized wrist anchors.
LagerNVS produces target wrist images at the head frame, tail frame, and selected keyframes.
These anchors are inserted into the target-view conditioning video at their true temporal slots.
The conditioned video is then encoded by the WAN inpainting VAE branch into a tensor $y$ containing mask channels and VAE latent channels.
This design keeps anchor evidence aligned with the temporal slots seen by the diffusion transformer; the stage-2 path does not directly overwrite target latent slots.

The current implementation supports three anchor regimes.
A head anchor supplies the initial wrist view.
A tail anchor constrains the far end of the clip and reduces long-horizon drift.
Keyframe anchors add sparse mid-trajectory evidence for contact and viewpoint changes.
During training, anchor dropout can remove evidence to improve robustness.
The visual ablation in Section~\ref{sec:experiments} evaluates these anchor choices, while downstream component ablations remain future work.

\subsection{External-View Temporal Memory}

The two external videos provide global workspace context and object motion.
Their VAE latents are selected as source latents and passed through a cross-view source projector to produce diffusion-transformer-dimensional source memory.
In temporal-local mode, each target latent frame attends to a local temporal window of source memory rather than the entire episode.
This constrains wrist synthesis with nearby external observations while reducing opportunities to attend to temporally inconsistent source evidence.

The source-memory branch is gated by a learned scalar or state-aware temporal gate.
The gate lets the model modulate how strongly external memory contributes during different phases of manipulation.
For example, source memory may be more helpful when the wrist view is heavily occluded, while anchor and action cues may dominate near contact.

\subsection{Action, Camera, and Scene Conditioning}

Robot state/action is embedded and injected into the latent stream.
The implementation supports a latent-time noise-injection mode and an AdaLN-style mode; the current cross-view experiments use the action/state sequence as a dynamic condition.
Target camera tokens can be projected by a time-wise MLP and added to the latent token tensor when enabled.

Scene tokens are built from cached tokens, camera-aware geometry sidecars, or runtime source frames.
In the current code path, these scene tokens are used by the 3D noise prior and optional alignment losses.
Although the DiT block contains a geometry-aware cross-attention helper, the scene-token pass-through to every DiT block is not active in the inspected implementation.
The paper therefore describes scene-token conditioning through the active sidecar/noise-prior path only.

\subsection{Action/View-Aware 3D Noise Initialization}

Standard diffusion starts from Gaussian noise independent of the scene, camera relation, and robot action phase.
We mix Gaussian noise $\epsilon$ with structured noise $\epsilon_{\mathrm{3D}}$ generated from scene tokens, source/target view IDs, spatial grid queries, and the action/state sequence:
\[
    \tilde{\epsilon} =
    \sqrt{1-\lambda_t^2}\,\epsilon + \lambda_t \epsilon_{\mathrm{3D}}.
\]
The mixture is normalized to match the scheduler's expected noise statistics.
At anchor slots, the effective mixing weight is attenuated using the inpainting mask from $y$, preventing the structured prior from overriding known anchor evidence.
The scheduler then builds the noisy latent and training target from $\tilde{\epsilon}$ rather than from independent Gaussian noise.

The prior has two active design choices.
First, it is view-aware: source and target view IDs and camera tokens condition the structured noise on the exo-to-wrist relation.
Second, it is action-aware: the action/state sequence modulates structured noise over latent time.
The prior remains noise, not a deterministic geometry reconstruction.
Its role is to bias early denoising toward a scene- and motion-consistent target view while preserving stochasticity for missing details.

\subsection{Training Objective}

For target latent $z_0$, timestep $t$, and mixed noise $\tilde{\epsilon}$, the scheduler constructs a noisy latent $z_t$ and a flow-matching training target $u_t$.
Stage-2 training minimizes mean-squared error over the target wrist latent sequence:
\[
    \mathcal{L}_{\mathrm{main}} =
    \left\| f_{\theta}(z_t, t, y, m_s, a, c) - u_t \right\|_2^2,
\]
where $m_s$ denotes source temporal memory.
Optional auxiliary losses include temporal consistency, target-state prediction, and scene-hidden alignment.
At inference, the same anchor, source-memory, action/state, camera, and 3D-noise paths condition wrist-view denoising.

\begin{algorithm}[t]
\caption{Guarded Stage-2 Training Step}
\label{alg:train}
\begin{algorithmic}[1]
\STATE Encode target wrist video into latent $z_0$.
\STATE Encode or load source-view latents and project them to temporal memory $m_s$.
\STATE Insert LagerNVS wrist anchors into the target conditioning video and encode inpainting condition $y$.
\STATE Embed action/state sequence and target camera tokens when available.
\STATE Build scene tokens from cached geometry sidecars or source frames.
\STATE Sample timestep $t$ and Gaussian noise $\epsilon$.
\STATE Generate structured noise $\epsilon_{\mathrm{3D}}$ and mix it with $\epsilon$ to obtain $\tilde{\epsilon}$.
\STATE Attenuate structured noise at known anchor slots and normalize mixed noise.
\STATE Construct noisy latent $z_t$ and scheduler target $u_t$.
\STATE Predict the target velocity/noise target with the conditional WAN denoiser.
\STATE Minimize the main MSE loss plus enabled auxiliary losses.
\end{algorithmic}
\end{algorithm}

\section{Experiments}
\label{sec:experiments}

The current experiments are a guarded evidence package rather than a final submission package.
All numbers in this section are copied from local CSV records in \texttt{experiments/}.
Missing analyses are marked explicitly.
\todo{PROTOCOL}: document one frozen split, resolution, frame count, inference budget, and input assumption set for all visual baselines.
\todo{CI}: report seeds, run counts, standard deviations, and confidence intervals for policy results.

\subsection{Datasets}

We use three data sources with complementary roles: DROID for real-world cross-view video generation, RLBench for controlled policy-grounded evaluation, and an in-house real-robot dataset for future physical validation.
The current paper reports DROID-style visual metrics and RLBench policy results; the in-house real-robot dataset is reserved for follow-up experiments and is not used for success-rate claims in this draft.

\textbf{DROID.}
DROID is a large-scale in-the-wild robot manipulation dataset collected with a standardized Franka Panda platform \cite{khazatsky2025droid}.
It contains approximately 76k successful trajectories, or 350 hours of interaction, across 564 scenes and diverse manipulation tasks.
Each episode records synchronized multi-view RGB-D observations from two external ZED 2 cameras and a wrist-mounted ZED Mini camera, together with robot states, action commands, natural-language annotations, and camera calibration.
This combination of external anchor views, calibrated geometry, and ground-truth wrist videos makes DROID a natural source for evaluating exo-to-wrist video generation under real-world scene variation.
In our experiments, DROID supports the visual generation protocol: the model receives two external-view videos with robot/action and camera information and predicts the corresponding wrist-view video.

\textbf{RLBench and PerAct-18.}
RLBench is a CoppeliaSim-based manipulation benchmark built around a Franka Panda arm and parallel-jaw gripper \cite{james2019rlbench}.
It provides language-conditioned task demonstrations with calibrated RGB-D observations from multiple viewpoints, including front, shoulder, and wrist cameras.
We use the 18-task evaluation suite popularized by PerAct \cite{shridhar2023peract}, which has become a common controlled benchmark for multi-task manipulation policies and recent VLA-style evaluations.
Compared with DROID, RLBench offers less visual diversity but tighter control over task definitions, camera geometry, and policy evaluation.
We therefore use RLBench to test whether generated wrist observations improve downstream policy training when real wrist views are withheld or scarce.

\subsection{Evaluation Questions}

The experiments ask three questions.
First, can the generator produce visually faithful wrist-view videos from two external views under the current DROID-style protocol?
Second, do the proposed conditioning components improve visual metrics in controlled ablations?
Third, do generated wrist videos help downstream policy training when real wrist observations are missing or incomplete?
The current draft answers these questions with existing visual metrics and RLBench summary results, but it does not claim DROID downstream control or real-robot deployment performance.

\subsection{Visual Generation Protocol}

This experiment tests whether a model can synthesize the held-out wrist-view video from two synchronized external-view videos, robot state/action, and camera information on DROID.
We first evaluate generic video fidelity with the standard reconstruction and distributional metrics used by video generation work: PSNR and SSIM measure pixel-level and structural agreement with the target wrist video; LPIPS measures perceptual frame similarity; FID measures frame-distribution realism; and FVD measures temporal video realism.
These metrics provide a necessary visual-quality audit because a wrist-view generator must preserve appearance, layout, and temporal coherence before it can be useful as a policy observation.
Table~\ref{tab:visual} reports this conventional comparison.
Baselines include wrist-view and exo-to-ego generation methods, a video diffusion baseline with wrist first-frame access, and geometry-only novel-view synthesis baselines.
Blank Wan2.2 and Cosmos rows from the source CSV are excluded because their metrics are unavailable.

\begin{table*}[t]
\caption{DROID-style visual wrist-view generation comparison from \texttt{experiments/Baseline.csv}. The result supports the guarded claim that \method{} improves over evaluated close baselines under the current protocol. \todo{PROTOCOL}}
\label{tab:visual}
\centering
\small
\begin{tabular}{llcccccc}
\toprule
Method & Type & Wrist first & PSNR $\uparrow$ & SSIM $\uparrow$ & LPIPS $\downarrow$ & FID $\downarrow$ & FVD $\downarrow$ \\
\midrule
VGGT projection \cite{wang2025vggt} & NVS & No & 8.93 & 0.2059 & 0.7843 & 208.98 & 834.02 \\
LagerNVS \cite{szymanowicz2026lagernvs} & NVS & No & 15.05 & 0.5130 & 0.4332 & 40.99 & 245.97 \\
SVD \cite{blattmann2023svd} & VDM & Yes & 12.84 & 0.4971 & 0.5294 & 31.87 & 363.84 \\
Exo2Ego-V \cite{liu2024exo2ego} & NVS+VDM & No & 15.43 & 0.4909 & 0.4020 & 22.73 & 130.11 \\
WristWorld \cite{qian2025wristworld} & NVS+VDM & No & 15.22 & 0.5161 & 0.3981 & 25.90 & 103.93 \\
\method{} & NVS+VDM & No & 19.46 & 0.6289 & 0.2174 & 14.39 & 57.92 \\
\bottomrule
\end{tabular}
\end{table*}

The conventional metrics show the strongest margins on temporal and perceptual fidelity: \method{} reduces FVD from 103.93 for WristWorld to 57.92 and reduces LPIPS from 0.3981 to 0.2174.
This result supports the statement that the proposed pipeline improves over evaluated close baselines in the current DROID protocol.
It should not yet be phrased as a universal state-of-the-art claim because baseline protocol, missing rows, and stochastic uncertainty still need audit.

Generic video metrics are not sufficient for robot manipulation videos.
A generated wrist video can be sharp and temporally smooth while still violating contact, gripper-object interaction, action following, or target-view geometry.
We therefore add a WorldArena-style embodied video evaluation \cite{shang2026worldarena}.
WorldArena was introduced to evaluate embodied world models across both perceptual and functional dimensions, and its analysis highlights a perception--functionality gap: high visual quality does not necessarily imply strong embodied utility.
This motivation directly matches wrist-view generation.
The generated video is not only a rendered observation; it should preserve contact timing, physically plausible object motion, perspective and depth cues from the target wrist camera, and controllability with respect to the robot action.

\begin{table*}[t]
\caption{WorldArena-style embodied video evaluation on DROID. Scores are normalized so higher is better. We report the aggregate EMScore and six perception sub-dimensions used by WorldArena: visual quality (VQ), motion quality (MQ), content consistency (CC), physics adherence (Phys.), 3D accuracy (3D), and controllability (Ctrl.).}
\label{tab:worldarena}
\centering
\scriptsize
\setlength{\tabcolsep}{4pt}
\begin{tabular}{llccccccc}
\toprule
Method & Type & EMScore $\uparrow$ & VQ $\uparrow$ & MQ $\uparrow$ & CC $\uparrow$ & Phys. $\uparrow$ & 3D $\uparrow$ & Ctrl. $\uparrow$ \\
\midrule
WristWorld \cite{qian2025wristworld} & NVS+VDM & 0.670 & 0.577 & 0.578 & 0.574 & 0.533 & 0.799 & 0.698 \\
Exo2Ego-V \cite{liu2024exo2ego} & NVS+VDM & 0.620 & 0.475 & 0.514 & 0.559 & 0.452 & 0.788 & 0.661 \\
Boundless World Model & VDM & 0.634 & 0.569 & 0.529 & 0.547 & 0.501 & 0.777 & 0.647 \\
SVD \cite{blattmann2023svd} & VDM & 0.624 & 0.518 & 0.447 & 0.563 & 0.535 & 0.780 & 0.698 \\
LagerNVS \cite{szymanowicz2026lagernvs} & NVS & 0.577 & 0.493 & 0.391 & 0.542 & 0.414 & 0.757 & 0.639 \\
VGGT projection \cite{wang2025vggt} & NVS & 0.640 & 0.513 & \textbf{0.848} & \textbf{0.806} & 0.211 & 0.678 & 0.505 \\
\method{} & NVS+VDM & \textbf{0.693} & \textbf{0.597} & 0.590 & 0.580 & \textbf{0.574} & \textbf{0.839} & \textbf{0.731} \\
\bottomrule
\end{tabular}
\end{table*}

Table~\ref{tab:worldarena} shows that \method{} obtains the highest aggregate EMScore among the evaluated methods on DROID, improving over WristWorld from 0.670 to 0.693.
The gains are concentrated in the dimensions most tied to robot usefulness: physics adherence, 3D accuracy, and controllability.
This pattern matters because some baselines score well on individual appearance-oriented axes without satisfying embodied constraints.
For example, VGGT projection has high motion quality and content consistency in this evaluation, but much lower physics adherence, 3D accuracy, and controllability, which is consistent with a geometry-heavy renderer that can preserve visible structure while missing contact- and action-conditioned dynamics.
Together, the conventional video metrics and WorldArena-style diagnostics support the intended claim: \method{} improves wrist-view video quality while better preserving the embodied cues required by downstream robot policies.

\subsection{Mechanism Ablations}

Table~\ref{tab:ablation} evaluates the visual effect of replacing placeholder conditioning with geometry anchors and then adding keyframe anchors, scene tokens, and the 3D noise prior.
The strongest current evidence is the progressive improvement across recorded rows.
The final 3D-noise variant improves over keyframe+scene from 19.27 to 19.46 PSNR, 0.6223 to 0.6397 SSIM, and 55.69 to 49.40 FVD.
This supports a visual-metric claim only; it does not by itself prove that the 3D noise prior causes downstream policy gains.

\begin{table*}[t]
\caption{Mechanism ablation from \texttt{experiments/Ablation.csv}. Missing FVD entries remain unavailable rather than silently filled. \todo{ABLATE}}
\label{tab:ablation}
\centering
\small
\begin{tabular}{lccccccccc}
\toprule
Variant & Head & Tail & Key & Scene & 3D & PSNR $\uparrow$ & SSIM $\uparrow$ & LPIPS $\downarrow$ & FVD $\downarrow$ \\
\midrule
Placeholder & No & No & No & No & No & 14.28 & 0.5042 & 0.4023 & -- \\
Head LagerNVS & Yes & No & No & No & No & 15.60 & 0.5244 & 0.3536 & -- \\
Head + scene & Yes & No & No & Yes & No & 16.09 & 0.5385 & 0.3531 & -- \\
Head + tail + scene & Yes & Yes & No & Yes & No & 18.13 & 0.5939 & 0.2637 & -- \\
Head + tail + keyframes & Yes & Yes & Yes & No & No & 18.51 & 0.6066 & 0.2360 & 62.36 \\
+ scene tokens & Yes & Yes & Yes & Yes & No & 19.27 & 0.6223 & 0.2219 & 55.69 \\
+ 3D-noise prior & Yes & Yes & Yes & Yes & Yes & 19.46 & 0.6397 & 0.2172 & 49.40 \\
\bottomrule
\end{tabular}
\end{table*}

\subsection{Policy-Grounded Evaluation on RLBench}

This experiment tests whether generated wrist-view videos are useful as policy observations rather than only visually plausible samples.
We evaluate on RLBench \cite{james2019rlbench}, using 18 manipulation tasks that span non-prehensile manipulation, pick-and-place, and high-precision insertion.
Each task provides expert demonstrations with language instructions, keyframes, and calibrated RGB-D observations from front, shoulder, and wrist cameras.
Policies are evaluated by task success rate; the current table reports mean success rates and still requires final seed, run-count, and confidence-interval reporting. \todo{CI}

We keep the RLBench task suite and evaluation protocol fixed, but restructure the training demonstrations into two disjoint subsets.
In the \emph{Available-Wrist Split}, real wrist observations are retained.
In the \emph{Wrist-Withheld Split}, wrist observations are removed from policy training and replaced with one of four alternatives: no wrist view, WristWorld-generated wrist videos, \method{}-generated wrist videos, or real wrist observations as an oracle reference.
This setup supports two policy-grounded questions.
The replacement setting asks whether generated wrist videos compensate for missing wrist cameras in the Wrist-Withheld Split.
The augmentation setting starts from the Available-Wrist Split and asks whether adding the Wrist-Withheld Split with synthetic wrist observations improves policy training.

\begin{figure*}[t]
\centering
\includegraphics[
    width=\textwidth,
    trim=115 100 105 80,
    clip
]{Figures/rlbench.pdf}
\caption{
Qualitative RLBench comparison.
Rows show external-view ground truth, wrist-view ground truth, \method{} generated wrist-view frames, and WristWorld-generated frames.
Compared with WristWorld, \method{} better preserves the target wrist-view geometry, object appearance, and temporal layout in these examples.
}
\label{fig:rlbench_qualitative}
\end{figure*}

Figure~\ref{fig:rlbench_qualitative} illustrates the visual difference in the downstream-control setting: the \method{} samples align more closely with the target wrist-view layout and object configuration.
Table~\ref{tab:rlbench_policy} then evaluates whether these generated observations improve policy learning.

\begin{table*}[t]
\centering
\scriptsize
\setlength{\tabcolsep}{3.6pt}
\caption{
Policy-grounded evaluation on RLBench.
Replacement evaluates alternatives for the Wrist-Withheld Split.
Augmentation starts from the Available-Wrist Split and adds the Wrist-Withheld Split with different wrist-view sources.
All values are average success rates in percentage; WW denotes WristWorld.
}
\label{tab:rlbench_policy}
\resizebox{\textwidth}{!}{%
\begin{tabular}{llccccccccccc}
\toprule
\multicolumn{2}{c}{Setting}
& \multicolumn{4}{c}{Replacement on Wrist-Withheld Split}
& \multicolumn{7}{c}{Augmentation from Available-Wrist Split} \\
\cmidrule(lr){1-2}
\cmidrule(lr){3-6}
\cmidrule(lr){7-13}
Dataset & Policy
& No Wrist
& WW
& Ours
& Real
& Avail. Only
& +WW
& +Ours
& Full Real
& Ours Gain
& Oracle Gap
& Recov. Gain \\
\midrule
RLBench-18 & RVT-2
& 70.78
& 68.91
& 72.13
& 72.49
& 70.58
& 76.00
& 78.98
& 81.87
& +8.40
& -2.89
& 74.4 \\
RLBench-18 & pi0.5
& 38.96
& 41.96
& 43.82
& 44.78
& 46.29
& 47.82
& 50.18
& 50.40
& +3.89
& -0.22
& 94.6 \\
RLBench-12 & pi0.5
& 57.43
& 61.97
& 64.53
& 66.50
& 68.03
& 69.10
& 73.57
& 73.73
& +5.53
& -0.17
& 97.1 \\
\bottomrule
\end{tabular}%
}
\vspace{2pt}

\footnotesize{
``Real'' and ``Full Real'' are oracle settings that use real wrist observations.
Ours Gain is ``+Ours'' minus ``Avail. Only''.
Oracle Gap is ``+Ours'' minus ``Full Real''.
Recov. Gain is the percentage of the full real-wrist augmentation gain recovered by synthetic wrist observations.
}
\end{table*}

In the replacement setting, \method{} improves over no-wrist training in all recorded policy/task groups and approaches the real-wrist oracle.
For RVT-2 on RLBench-18, generated wrist observations reach 72.13\% average success, 0.36 points below real wrist observations and 3.22 points above WristWorld.
For pi0.5, \method{} raises RLBench-18 success from 38.96\% to 43.82\%, compared with 44.78\% for real wrist observations.
On the RLBench-12 subset, \method{} raises pi0.5 success from 57.43\% to 64.53\%, compared with 66.50\% for the real-wrist reference.
These results support the bounded claim that generated wrist observations recover much of the downstream value of real wrist cameras in the recorded RLBench settings.

The augmentation results show a stronger policy-training signal.
Starting from the Available-Wrist Split, adding \method{} generated wrist observations from the Wrist-Withheld Split improves every recorded group.
For RVT-2, \method{} improves success from 70.58\% to 78.98\%, outperforming WristWorld augmentation by 2.98 points and recovering 74.4\% of the full real-wrist augmentation gain.
For pi0.5, synthetic augmentation nearly matches the full real-wrist oracle in these records: 50.18\% versus 50.40\% on RLBench-18, and 73.57\% versus 73.73\% on RLBench-12.
We therefore phrase the policy claim as supplementing and approaching real wrist observations, not as universally replacing them across policies and tasks.

\begin{table*}[t]
\centering
\tiny
\setlength{\tabcolsep}{2.1pt}
\caption{
Detailed per-task success rates on RLBench-18.
The table expands the aggregate results in Table~\ref{tab:rlbench_policy}.
All values are success rates in percentage; WW denotes WristWorld.
}
\label{tab:rlbench_18_full}
\resizebox{\textwidth}{!}{%
\begin{tabular}{ccl*{19}{c}}
\toprule
Dataset & Policy & Training Setting
& Avg.
& PID & RAD & TT & SBCT & OD & PGC & PSS & PMS & PB & CJ
& SB & PC & PWR & LBI & STDS & ISP & MOG & SC \\
\midrule
RLBench-18 & RVT-2 & Withheld Real
& 72.49
& 97.6 & 99.2 & 96.0 & 61.6 & 76.8 & 50.4 & 47.2 & 89.6 & 87.2 & 91.6
& 59.6 & 15.6 & 93.2 & 83.6 & 98.8 & 26.8 & 97.2 & 32.8 \\
& & Withheld Ours
& 72.13
& 95.2 & 100.0 & 90.8 & 72.4 & 77.2 & 40.0 & 48.8 & 89.2 & 86.8 & 97.6
& 55.2 & 31.2 & 92.0 & 81.2 & 100.0 & 19.6 & 99.2 & 22.0 \\
& & Withheld No Wrist
& 70.78
& 97.6 & 98.0 & 90.4 & 60.8 & 73.6 & 42.8 & 38.8 & 89.2 & 96.0 & 88.0
& 54.4 & 18.4 & 92.0 & 71.6 & 94.8 & 25.6 & 99.2 & 42.8 \\
& & Withheld WW
& 68.91
& 86.8 & 95.2 & 91.6 & 58.8 & 83.2 & 38.0 & 38.8 & 81.2 & 94.4 & 88.0
& 48.8 & 14.8 & 93.2 & 63.6 & 100.0 & 34.4 & 98.0 & 31.6 \\
& & Avail. Only
& 70.58
& 99.6 & 99.6 & 87.2 & 88.4 & 78.0 & 41.6 & 32.8 & 82.8 & 83.6 & 80.4
& 56.0 & 26.0 & 89.6 & 83.6 & 95.2 & 23.6 & 99.2 & 23.2 \\
& & Avail. + Ours
& 78.98
& 98.4 & 99.6 & 88.8 & 92.4 & 78.8 & 63.6 & 37.6 & 94.4 & 67.6 & 95.6
& 63.2 & 48.0 & 95.6 & 90.4 & 93.2 & 24.8 & 99.6 & 90.0 \\
& & Avail. + WW
& 76.00
& 98.8 & 94.4 & 91.2 & 65.6 & 74.8 & 64.0 & 43.2 & 90.0 & 95.6 & 96.4
& 39.2 & 22.8 & 93.2 & 89.6 & 100.0 & 42.0 & 98.0 & 69.2 \\
& & Full Real Oracle
& 81.87
& 97.6 & 98.0 & 91.2 & 95.2 & 79.2 & 78.8 & 47.6 & 92.4 & 88.4 & 99.6
& 73.2 & 34.0 & 92.8 & 88.8 & 100.0 & 36.4 & 98.4 & 82.0 \\
\cmidrule(lr){2-22}
RLBench-18 & pi0.5 & Withheld Real
& 44.78
& 32.4 & 91.6 & 57.6 & 95.6 & 33.6 & 25.2 & 5.6 & 84.8 & 56.4 & 94.4
& 0.4 & 0.0 & 47.2 & 1.2 & 79.2 & 0.8 & 100.0 & 0.0 \\
& & Withheld Ours
& 43.82
& 46.0 & 93.2 & 53.2 & 78.0 & 32.0 & 35.2 & 6.0 & 73.2 & 63.2 & 92.4
& 3.2 & 0.0 & 54.0 & 4.8 & 59.2 & 0.0 & 94.8 & 0.4 \\
& & Withheld No Wrist
& 38.96
& 33.6 & 92.4 & 48.8 & 86.4 & 56.8 & 7.2 & 4.4 & 79.6 & 44.8 & 70.4
& 0.8 & 0.4 & 48.4 & 6.0 & 26.8 & 0.0 & 94.0 & 0.4 \\
& & Withheld WW
& 41.96
& 40.0 & 92.0 & 44.0 & 76.8 & 32.4 & 34.8 & 4.0 & 65.6 & 60.4 & 91.2
& 2.0 & 0.0 & 55.2 & 5.6 & 57.2 & 0.0 & 94.0 & 0.0 \\
& & Avail. Only
& 46.29
& 51.6 & 96.0 & 64.0 & 87.6 & 37.2 & 16.4 & 7.6 & 82.8 & 67.6 & 81.2
& 0.4 & 0.4 & 53.6 & 6.4 & 82.0 & 2.0 & 96.4 & 0.0 \\
& & Avail. + Ours
& 50.18
& 47.2 & 96.4 & 56.0 & 95.2 & 24.4 & 39.2 & 5.6 & 96.8 & 98.8 & 100.0
& 2.0 & 0.0 & 46.4 & 8.0 & 82.8 & 1.6 & 99.6 & 3.2 \\
& & Full Real Oracle
& 50.40
& 46.0 & 91.6 & 66.0 & 97.2 & 27.6 & 47.6 & 5.2 & 95.2 & 94.8 & 100.0
& 4.4 & 0.0 & 31.6 & 5.2 & 88.8 & 1.2 & 98.4 & 6.4 \\
\bottomrule
\end{tabular}%
}
\vspace{2pt}

\footnotesize{
Task abbreviations:
PID: put item in drawer;
RAD: reach and drag;
TT: turn tap;
SBCT: slide block to color target;
OD: open drawer;
PGC: put groceries in cupboard;
PSS: place shape in shape sorter;
PMS: put money in safe;
PB: push buttons;
CJ: close jar;
SB: stack blocks;
PC: place cups;
PWR: place wine at rack location;
LBI: light bulb in;
STDS: sweep to dustpan of size;
ISP: insert onto square peg;
MOG: meat off grill;
SC: stack cups.
}
\end{table*}

The task-level results in Table~\ref{tab:rlbench_18_full} show where synthetic wrist observations help and where gaps remain.
\method{} is particularly useful for several manipulation-heavy tasks, while the remaining gap to the full real-wrist oracle is concentrated in fine-grained spatial and contact-rich tasks such as insertion, stacking, and multi-object placement.
This pattern is consistent with the aggregate result: generated wrist videos recover substantial downstream utility, but high-precision contact remains the hardest regime.

\subsection{Missing Analyses}

The current draft needs three analyses before final claims.
First, downstream component ablations should train policies with videos from at least the full model, no keyframe anchors, no scene tokens, and no 3D noise prior.
Second, robustness splits should isolate fast wrist motion, severe occlusion, small-object contact, and unseen backgrounds.
Third, policy-aware diagnostics such as contact-phase accuracy, gripper/object tracking consistency, VLA feature distance, or action prediction error would connect visual metrics to control behavior.

\section{Limitations and Scope}

This draft intentionally keeps claims narrow.
The current evidence supports DROID-style visual generation and RLBench downstream policy results, but does not support DROID downstream policy success or real-robot success.
The visual protocol and ablation tables need freezing before final submission.
RLBench policy summaries need statistical reporting.
The current code map also shows that per-block scene-token cross-attention exists as a helper but is not active in the inspected \texttt{model\_fn\_wan\_video} path.
Any final manuscript text that claims active DiT scene-token cross-attention must be revised unless that branch is enabled and rerun.

The method also inherits limitations from the anchor source and target domain.
If LagerNVS anchors are geometrically wrong, the diffusion model may receive misleading target-view evidence.
If the two external cameras do not observe the contact region, external memory cannot fully recover hidden object state.
If the downstream policy relies on tactile or force cues not visible in wrist RGB, a plausible generated wrist video may still be insufficient.
These limitations motivate hard-case splits and policy-aware diagnostics rather than only reporting average visual metrics.

\section{Conclusion}

We presented \method, a geometry-anchor guided wrist-view video generator for robot manipulation.
The core idea is to use sparse target-view geometry as slot-aligned evidence for video diffusion, complemented by external-view temporal memory, action/state conditioning, and an action/view-aware 3D noise prior.
Current results indicate improved visual metrics under the evaluated DROID-style protocol and policy benefits on RLBench.
Before submission, the paper must close the novelty, protocol, statistical, downstream-ablation, figure, and code-claim audit gates marked throughout the draft.

\bibliography{references}

% Uncomment only after final AAAI-27 instructions confirm checklist placement.
% \input{ReproducibilityChecklist.tex}

\end{document}
